\documentclass[10pt,a4paper]{article}
\usepackage[margin=2.5cm]{geometry}
\usepackage[utf8]{inputenc}
\usepackage[T1]{fontenc}
\usepackage[spanish,es-nodecimaldot]{babel}
\usepackage{amsmath,amssymb,booktabs,array}
\usepackage[colorlinks=true,urlcolor=blue]{hyperref}
\title{Diagnóstico estructural del escalado\\en redes multicommodity}
\author{Nota separada del manuscrito OMS}
\date{18 de octubre de 2026}
\begin{document}\maketitle
\noindent\textbf{Resultado.} Se encontraron modelos externos con bloques locales,
acoplamiento efectivo y contratos activos. Las 773 filas de acoplamiento de doce
modelos cambian tanto con Role-Hybrid como al imponer el presupuesto $10^{-5}$.
Sin embargo, las 210 soluciones de siete políticas ya satisfacen $10^{-7}$,
y la réplica no acredita una ventaja general frente al modelo sin escalar.
El manuscrito OMS v1.3.14 permanece intacto.

\section{Selección por estructura antes de resolver}
La pregunta era adecuada: la selección debe permitir probar el mecanismo que
motiva el método. Se fijaron seis problemas R de cargos fijos de CommaLAB,
R01.1 a R06.1, y las dos primeras matrices de demandas por orden lexicográfico
de cada uno de los tres grupos Günlük de expansión de capacidad.\footnote{\url{https://commalab.di.unipi.it/datasets/mmcf/}}
Los problemas R tienen diez nodos, 35 o 60 arcos y 10, 25 o 50 mercancías.
Los grupos Günlük tienen respectivamente 15/22/210, 27/51/702 y 16/49/240
nodos/enlaces/mercancías según los archivos nativos. Las dos demandas de cada
grupo comparten matriz: no son seis matrices independientes.
Las dos colecciones constituyen formulaciones y procedencias distintas dentro
de redes; no equivalen a dos dominios de aplicación independientes.

La formulación R conserva balances $Bx_k=d_k$ y capacidades
$\sum_k x_{ke}\le U_e y_e$, con $y_e$ binaria.
La expansión utiliza capacidad bidireccional: cada sentido queda acotado por
la misma capacidad instalada $u_e+y_{e,1}+\lambda y_{e,2}$, con instalaciones
enteras no negativas. No se suman ambos sentidos como en capacidad no dirigida.
La distinción está respaldada por la clasificación de los autores.\footnote{\url{https://arxiv.org/html/1711.10147}, secciones 1 y 3.}
Los datos nativos no se perturban ni se reescalan artificialmente.
El rango de coeficientes es 613--2640 en R y solo 4 en Günlük: allí la
densidad de las filas activa L2 sin un rango extremo de coeficientes.

También se inspeccionaron los diez problemas presentes en el ZIP SSLP de
SIPLIB.\footnote{\url{https://www2.isye.gatech.edu/~sahmed/siplib/sslp/sslp.html}}
Sus filas de recurso tocan un único escenario y las columnas de primera etapa:
no hay filas entre dos bloques de recurso distintos. No se introdujeron copias
de variables para forzar una estructura diferente.

\section{Qué hipótesis se cumplen}
Los doce modelos tienen la estructura angular
\[
 A=\begin{bmatrix}I_K\otimes B&0\\
 \mathbf{1}_K^{\mathsf T}\otimes I&D\end{bmatrix},
\]
donde la última fila de bloques representa varias filas de capacidad.
Los bloques $B$ son incidencias rectangulares, con rango $n-1$ y dependencias
de conservación. Por tanto, \textbf{no se aplica el teorema triangular con bloques
identidad}: escalar filas no convierte estas incidencias en identidades.
Los contratos de residuos sí son aplicables sin esa hipótesis.

Todos los coeficientes locales no nulos son $\pm1$. La media geométrica local
elige uno y, en estos modelos, Role-Hybrid coincide con Coupling-only.
La etapa de acoplamiento actúa sobre las 773 filas, a diferencia de la ablación
histórica de despacho. Esto permite aislar el mecanismo, sin atribuirle
información adicional de propiedad de columnas.

El diagnóstico conserva todas las filas. Calcula $\kappa_2^+$ mediante una
descomposición ortogonal analítica entre el promedio de las mercancías y sus
contrastes; no transforma las variables del MIP. La fórmula se contrastó con
120 SVD completas y la equilibración se contrastó elemento a elemento.
Son medidas de la matriz completa, no del condicionamiento de una base LP.
Las diez pasadas diádicas de equilibración constituyen un sustituto declarado;
no observan la implementación interna propietaria del solver.

\begin{center}\small
\begin{tabular}{lrrrr}
\toprule
&\multicolumn{2}{c}{R01.1}&\multicolumn{2}{c}{Noruega, demanda 1}\\
Política&Antes&Tras 10 pasadas&Antes&Tras 10 pasadas\\\midrule
Base&320.44&3.39&35.89&18.32\\
Role-Hybrid&3.42&3.42&5.50&5.95\\
Flat-L2&1.58&2.51&7.87&5.79\\
Budget-L2&82.17&2.46&11.48&10.32\\\bottomrule
\end{tabular}
\end{center}
La mejora inicial de centrado o condición no implica superioridad posterior.
Las doce instancias mejoran con Role-Hybrid antes de equilibrar; después,
solo seis superan a Base según esta medida. Flat-L2 supera a Base en las doce.

\section{Contratos activos y verificación independiente}
Se fijaron metas absolutas $10^{-3}$, $10^{-5}$ y $10^{-7}$ en unidades de flujo
de la fuente, como objetivos del estudio. La hipótesis de residuo exportado es
$\epsilon=10^{-6}$ y la distancia uniforme de redondeo entero es
$\eta=10^{-9}$. Los factores usan ventana $[10^{-9},10^9]$ y banda cerrada
de identidad $[1/2,2]$. Son adaptadores Python explícitos; no se afirma identidad
binaria con el exportador C++.

\begin{center}\small
\begin{tabular}{lrrr}\toprule
Presupuesto&Filas de acoplamiento&Factores modificados&Incompatibles\\\midrule
$10^{-3}$&773&40&0\\
$10^{-5}$&773&773&0\\
$10^{-7}$&773&517&256\\\bottomrule
\end{tabular}
\end{center}
Los contratos incompatibles corresponden a la cota uniforme de redondeo en
problemas R. No prueban inviabilidad del modelo ni incumplimiento del presupuesto
por soluciones observadas. En una capacidad R, el presupuesto suficiente sin
defecto de exportación es
\[
 \delta\ge\epsilon/\alpha+U_e\eta.
\]
El término $U_e\eta$ no se corrige mediante escalado de filas. Si las variables
enteras entregadas son exactamente enteras, su error real de redondeo es cero:
la incompatibilidad uniforme puede coexistir con una solución precisa.

Las 60 exportaciones se comprobaron tras importarlas: matriz, RHS, nombres,
objetivo, cotas y dominios. Las multiplicaciones por factores reales redondean
los datos; el verificador mide exactamente el defecto puntual. Los balances,
capacidades, cotas, dominios y costos se reconstruyen desde arcos y demandas
nativos con racionales, separando punto bruto y redondeado. No se usan los
atributos de factibilidad del optimizador. Se verifican las hipótesis de
transferencia, incluido el defecto de exportación; el parámetro nominal de
tolerancia no sustituye esas comprobaciones.

\section{Piloto y réplica con calidad común}
El piloto conserva 84 ejecuciones: doce modelos y siete políticas, con semilla
223, Gurobi 13.0.2, cuatro hilos, presolve y escalado internos automáticos,
WorkLimit$=0.25$, TimeLimit$=10$ s y brecha nominal del 1\%.
IntFeasTol$=10^{-9}$ es común. Se comparan Base, Ruiz-R4, Flat-L2,
Role-Hybrid, Budget-L2 a $10^{-5}$, Flat-L2-tight con tolerancia primal
$10^{-9}$ y Flat-L2 con NumericFocus$=3$.
Coupling-only es un alias demostrado y no se cuenta como evidencia adicional.
Los seis problemas de expansión alcanzan el límite de trabajo: sus incumbentes
no se presentan como óptimos certificados.

Las 84 soluciones cumplen $10^{-7}$ con todos los controles. El mayor residuo
original bruto o redondeado es $2.3284\cdot10^{-12}$. Budget-L2 usa menos
trabajo que L2 en las seis R, pero los costos y los esfuerzos de cierre de cota
difieren. Esa comparación inicial no basta para declarar una ventaja práctica.

Se registró después una réplica sobre las seis R completas, con semillas nuevas
227, 229 y 233: 126 ejecuciones. Todas las políticas buscan el mismo costo,
1.01 veces un costo factible exacto del piloto, mediante BestObjStop, conservando
los otros límites. Es una referencia factible, no un óptimo certificado.
Las 126 entregas cumplen el objetivo común y $10^{-7}$; el mayor residuo es
$1.2435\cdot10^{-14}$.

\begin{center}\small
\begin{tabular}{lrr}\toprule
Control&Trabajo Budget/control&Ciclo Budget/control\\\midrule
Base&1.009&1.129\\
Flat-L2&0.884&1.071\\
Flat-L2-tight&0.884&1.067\\
Flat-L2-NF3&0.599&0.970\\
Role-Hybrid&0.905&1.089\\
Ruiz-R4&0.921&0.921\\\bottomrule
\end{tabular}
\end{center}
Se informan medias geométricas de 18 razones emparejadas. El ciclo suma
construcción, selección y exportación medidas por separado, carga y resolución;
excluye verificación posterior. Los procesos son nuevos y secuenciales. Los
tiempos son cortos, descriptivos y sensibles a la máquina compartida; las
semillas no convierten seis modelos en dieciocho instancias independientes.
El ciclo parte de los datos nativos cargados y excluye arranque del intérprete,
lectura inicial e inicialización de la licencia.

\section{Decisión científica}
La selección estructural fue pertinente: se comprobaron acoplamiento activo,
contratos activos y cambios que sobreviven parcialmente a la equilibración
de diagnóstico. No se encontró una recuperación exclusiva de precisión;
tampoco una ventaja general sobre Base a calidad común. La reducción de
trabajo frente a L2 no se traduce en menor ciclo total medido.
Estos resultados se conservan fuera del artículo: no justifican presentar
una nueva contribución competitiva ni extender el alcance del teorema.
Buscar recuperación de precisión exige además dificultades numéricas nativas:
un contrato activo no demuestra que la solución sin contrato viole el presupuesto.

\noindent\textbf{Reproducción.} La carpeta incluye archivos nativos, diseños,
60 exportaciones, 210 registros del optimizador, puntos, comprobaciones racionales
y sumarios. Los verificadores funcionan sin licencia comercial. Las URL y
SHA-256 quedan vinculadas en los manifiestos; se conserva la evidencia adversa.
\end{document}
